\section{Reasoning scaffold: hacer que el proceso intermedio pueda buscarse y revisarse}

La última parte de contenido sobre modelos del curso trata Chain-of-Thought, Tree of Thoughts y Socratic Questioning. La idea común es cambiar «una salida directa» por un proceso intermedio de varios pasos, de modo que el modelo obtenga más test-time compute y el sistema conserve puntos donde revisar y retroceder.

\subsection{Por qué CoT puede ser útil}

Chain-of-Thought (cadena de pensamiento, CoT) pide al modelo, mediante ejemplos o instrucciones, que genere pasos intermedios. En tareas aritméticas y simbólicas de varios pasos es difícil asignar directamente la respuesta a la pregunta; al descomponer la tarea, la relación local de cada paso puede quedar más cerca de la distribución de entrenamiento. Las dos páginas siguientes muestran primero la intermediate guidance y después comparan la diferencia de información entre un standard prompt y un CoT example.

\lecturefigure{slide-070.jpg}{La intermediate guidance despliega una sola predicción en un razonamiento de varios pasos}{Slide deck oficial de CS25 V4, página 70}

\lecturefigure{slide-071.jpg}{Comparación entre prompting estándar y Chain-of-Thought prompting}{Slide deck oficial de CS25 V4, página 71}

\begin{importantbox}{Descomposición probabilística de CoT}
Si la reasoning trace intermedia es $r$ y la respuesta final es $y$, el proceso de generación puede escribirse como
\[
p(y\mid x)=\sum_r p(r\mid x)p(y\mid x,r).
\]
En la práctica, el decoding solo muestrea o busca unos cuantos $r$. La ventaja de CoT proviene de un cómputo intermedio y una descomposición más ricos; no garantiza que la trace muestreada refleje fielmente el proceso causal interno del modelo.
\end{importantbox}

\teachervoice{El ponente usa la multiplicación de números de diez cifras para explicar la intuición de los pasos intermedios: las personas también suelen calcular por pasos. El valor adicional de CoT es que los errores quedan visibles: el lector puede localizar en qué paso el modelo omitió o malinterpretó algo, en lugar de ver solo una respuesta incorrecta.}

\subsection{Más allá de la mejora de desempeño: hay que ver el tipo de error}

CoT suele producir mejoras en modelos grandes, pero todavía genera pasos irrelevantes, errores locales o conclusiones erróneas tras un proceso correcto. La clase divide los errores en missing step y semantic misunderstanding, y destaca en particular que los modelos pequeños pueden ser inestables incluso en el mapping local. En la lectura de las figuras de esta sección, lo central es la error taxonomy y el model size, no solo recordar que el desempeño promedio sube.

\lecturefigure{slide-072.jpg}{Resumen de la clase sobre la mejora de CoT y los errores restantes}{Slide deck oficial de CS25 V4, página 72}

\begin{warningbox}{Un rationale largo no equivale a un reasoning confiable}
El modelo puede generar explicaciones con formato completo pero sin relación con la respuesta, y también puede redactar la justificación hacia atrás después de obtener la respuesta. La validación debe examinar si los estados intermedios son ejecutables y verificables y si admiten un counterfactual change, no solo si el texto parece un razonamiento.
\end{warningbox}

\lecturefigure{slide-073.jpg}{Errores de missing-step y de semantic-understanding en modelos pequeños}{Slide deck oficial de CS25 V4, página 73}

\begin{knowledgebox}{Por qué los modelos pequeños fallan con más facilidad en CoT}
CoT aumenta la longitud de la secuencia y la coherencia que debe mantenerse. Si el modelo carece de capacidades semánticas o aritméticas básicas, más tokens crean más puntos donde equivocarse. Las rutas de mejora incluyen supervision por pasos, tool execution, verifier, distillation y una state representation más estructurada.
\end{knowledgebox}

\subsection{Generalized CoT: el razonamiento no tiene una sola línea}

El curso plantea que la definición de CoT es demasiado rígida: distintos problemas pueden requerir búsqueda con ramificaciones, retroceso, recursión sobre subproblemas o herramientas externas. Lo central del generalized reasoning es elegir el computation graph adecuado. Las tres páginas siguientes pasan de una trace lineal a la tree search y la recursive decomposition, y comparan state, transition y verifier.

\lecturefigure{slide-074.jpg}{De un Chain-of-Thought fijo hacia estructuras de razonamiento más generales}{Slide deck oficial de CS25 V4, página 74}

\begin{importantbox}{Tres ejes de diseño de un reasoning scaffold}
\begin{enumerate}
  \item \textbf{State}: si el estado intermedio es lenguaje natural, un programa, una ecuación o una observation del entorno;
  \item \textbf{Transition}: si el siguiente paso proviene de un muestreo del modelo, de reglas, de una búsqueda o de la ejecución de una herramienta;
  \item \textbf{Verifier}: cómo se puntúa, poda, retrocede y termina cada estado local.
\end{enumerate}
Solo cuando los tres quedan definidos con claridad, un método de razonamiento deja de ser un simple estilo de prompt.
\end{importantbox}

\subsection{Tree of Thoughts y Socratic decomposition}

Tree of Thoughts (árbol de pensamientos, ToT) mantiene varios estados candidatos y realiza lookahead/backtracking; Socratic Questioning (cuestionamiento socrático) plantea subproblemas de forma recursiva y luego combina sus respuestas en el problema original. Ambos intercambian más inference compute por espacio de búsqueda.

\lecturefigure{slide-075.jpg}{Ramificación, evaluación y retroceso en Tree of Thoughts}{Slide deck oficial de CS25 V4, página 75}

\begin{importantbox}{Búsqueda de anchura limitada}
Si cada estado se expande en $b$ candidatos, se conservan los top-$k$ y la profundidad es $d$, el número de llamadas al modelo crece aproximadamente como $O(dkb)$. El sesgo del verifier determina qué ramas se eliminan prematuramente; por ello, la calidad de la búsqueda no depende solo del generator, sino también del state scoring.
\end{importantbox}

\lecturefigure{slide-076.jpg}{Socratic Questioning genera y responde subproblemas de forma recursiva}{Slide deck oficial de CS25 V4, página 76}

\begin{knowledgebox}{Lectura de la figura: diferencia entre decomposition y search}
La decomposition divide un problema grande en problemas pequeños relativamente independientes; la search elige entre varias rutas candidatas. El Socratic method pone más énfasis en la estructura de subproblemas; ToT, en los estados candidatos y el retroceso. Los sistemas complejos suelen combinar ambos y llamar en los nodos hoja a un calculator, un code executor o a retrieval.
\end{knowledgebox}

\subsection{Resumen de la sección}

El valor de un reasoning scaffold está en aumentar el cómputo intermedio, exponer los errores y ofrecer interfaces de búsqueda y verificación. CoT, ToT y Socratic decomposition no producen automáticamente un reasoning genuino; solo se acercan a un proceso de cómputo controlable cuando se combinan con estados estructurados, tools, verifier y stopping rule.
